\documentclass{article}
\usepackage{amsmath,amssymb,amsthm}
\usepackage{algorithm}
\usepackage{algorithmic}
\usepackage{enumitem}
\usepackage{hyperref}
\usepackage{geometry}
\geometry{margin=1in}
\newtheorem{theorem}{Theorem}
\newtheorem{lemma}{Lemma}
\newtheorem{assumption}{Assumption}
\newtheorem{corollary}{Corollary}
\newcommand{\E}{\mathbb{E}}
\newcommand{\norm}[1]{\left\lVert #1\right\rVert}
\newcommand{\ip}[2]{\langle #1,#2\rangle}
\newcommand{\R}{\mathbb{R}}

\begin{document}

\section*{Algorithm Name}
\textbf{Frank–Wolfe (Conditional Gradient) Method for Non‑Convex Smooth Optimization}

\section*{Mathematical Formulation}
Let $f:\R^{d}\rightarrow\R$ be differentiable and let $\mathcal C\subset\R^{d}$ be a non‑empty compact convex set.  
At iteration $k\ge 0$ we perform

\begin{align}
s_k &:= \arg\min_{s\in\mathcal C}\;\ip{\nabla f(x_k)}{s}, \tag{LMO}\\[4pt]
\gamma_k &\in (0,1] \quad\text{(step‑size, possibly diminishing)}\\[4pt]
x_{k+1} &:= (1-\gamma_k)\,x_k + \gamma_k\, s_k . \tag{FW\;update}
\end{align}

The quantity 
\[
g_{\mathcal C}(x_k):=\max_{s\in\mathcal C}\ip{-\nabla f(x_k)}{s-x_k}
       =\ip{-\nabla f(x_k)}{s_k-x_k}
\]
is called the \emph{Frank–Wolfe gap} and serves as a stationarity measure.

\section*{Pseudocode}
\begin{algorithm}[H]
\caption{Frank–Wolfe for Non‑Convex $f$}
\begin{algorithmic}[1]
\Require Initial point $x_0\in\mathcal C$, step‑size rule $\{\gamma_k\}_{k\ge0}$.
\For{$k=0,1,2,\dots$}
    \State Compute gradient $g_k\gets\nabla f(x_k)$.
    \State \textbf{LMO:} $s_k \gets \displaystyle\arg\min_{s\in\mathcal C}\ip{g_k}{s}$.
    \State Choose step size $\gamma_k\in(0,1]$ (e.g. $\gamma_k=\frac{2}{k+2}$ or line‑search).
    \State Update $x_{k+1}\gets (1-\gamma_k)x_k+\gamma_k s_k$.
    \If{termination criterion satisfied (e.g. $g_{\mathcal C}(x_k)\le\varepsilon$)}
        \State \textbf{break}
    \EndIf
\EndFor
\State \Return $x_{k}$.
\end{algorithmic}
\end{algorithm}

\section*{Dry Run Example}
Consider the scalar problem  

\[
f(w) = (w-3)^2,\qquad\mathcal C=[0,5],\qquad w_0=0,\qquad \gamma_k\equiv\gamma=0.1 .
\]

\begin{center}
\begin{tabular}{c|c|c|c|c}
$k$ & $w_k$ & $\nabla f(w_k)=2(w_k-3)$ & $s_k$ (LMO) & $w_{k+1}=(1-\gamma)w_k+\gamma s_k$ \\ \hline
0 & $0$ & $-6$ & $\displaystyle\arg\min_{s\in[0,5]}(-6)s =5$ & $0.9\cdot0+0.1\cdot5 = 0.5$ \\[4pt]
1 & $0.5$ & $-5$ & $5$ & $0.9\cdot0.5+0.1\cdot5 = 0.95$ \\[4pt]
2 & $0.95$ & $-4.1$ & $5$ & $0.9\cdot0.95+0.1\cdot5 = 1.355$ \\[4pt]
3 & $1.355$ & $-3.29$ & $5$ & $0.9\cdot1.355+0.1\cdot5 = 1.7195$
\end{tabular}
\end{center}

Objective values: $f(0)=9$, $f(0.5)=6.25$, $f(0.95)=4.2025$, $f(1.355)=2.711$, $f(1.7195)=1.633$.
The Frank–Wolfe gap $g_{\mathcal C}(w_k)=\ip{-\nabla f(w_k)}{s_k-w_k}=(-\nabla f(w_k))(5-w_k)$ decreases accordingly.

\section*{Assumptions}
\begin{assumption}[Smoothness]\label{ass:smooth}
$f$ is $L$‑smooth on an open set containing $\mathcal C$; i.e.,
\[
\forall x,y\in\mathcal C:\qquad 
f(y)\le f(x)+\ip{\nabla f(x)}{y-x}+\frac{L}{2}\norm{y-x}^2 .
\]
\end{assumption}

\begin{assumption}[Compactness]\label{ass:compact}
$\mathcal C$ is compact and has bounded diameter
\[
D := \max_{x,y\in\mathcal C}\norm{x-y}<\infty .
\]
\end{assumption}

\begin{assumption}[Step‑size]\label{ass:stepsize}
The step size satisfies $0<\gamma_k\le 1$ and
\[
\sum_{k=0}^{\infty}\gamma_k = \infty,\qquad 
\sum_{k=0}^{\infty}\gamma_k^2 < \infty .
\]
A concrete choice is $\gamma_k = \frac{2}{k+2}$.
\end{assumption}

\section*{Main Convergence Theorem}
\begin{theorem}[Convergence of Frank–Wolfe for Non‑Convex Smooth $f$]\label{thm:main}
Under Assumptions \ref{ass:smooth}–\ref{ass:stepsize}, let $\{x_k\}$ be the sequence generated by the algorithm. Define the minimal Frank–Wolfe gap up to iteration $T$ as
\[
\tilde g_T := \min_{0\le k\le T} g_{\mathcal C}(x_k).
\]
Then for any $T\ge 0$,
\[
\tilde g_T \;\le\; \frac{2\bigl(f(x_0)-f_{\inf}\bigr)}{\sum_{k=0}^{T}\gamma_k}
               \;+\; L D^2 \frac{\sum_{k=0}^{T}\gamma_k^2}{\sum_{k=0}^{T}\gamma_k},
\tag{1}
\]
where $f_{\inf}:=\inf_{x\in\mathcal C}f(x)$. In particular, for the canonical step‑size $\gamma_k=2/(k+2)$,
\[
\tilde g_T = O\!\left(\frac{1}{\sqrt{T}}\right).
\]
\end{theorem}

\section*{Theoretical Properties}
\begin{itemize}
\item \textbf{Stability.} The iterates remain in $\mathcal C$ for all $k$ because the update is a convex combination of two points in $\mathcal C$.
\item \textbf{Hyper‑parameter sensitivity.} The convergence bound (1) depends on the ratio
\(\displaystyle\frac{\sum_{k=0}^{T}\gamma_k^2}{\sum_{k=0}^{T}\gamma_k}\).
Choosing $\gamma_k\asymp 1/k$ makes both numerator and denominator of order $\log T$, yielding the $O(1/\sqrt{T})$ rate. Larger constant step sizes accelerate early progress but increase the $L D^2$ term.
\item \textbf{Comparison to baselines.} For convex $f$, the same analysis yields the optimal $O(1/T)$ rate for the dual gap; in the non‑convex case we recover the best known $O(1/\sqrt{T})$ rate for first‑order methods that only access a linear minimization oracle.
\end{itemize}

\section*{Discussion}
The theorem shows that even without convexity the Frank–Wolfe gap—an intrinsic stationarity measure—vanishes at the same sub‑linear speed as stochastic gradient methods. The proof hinges only on smoothness and bounded domain, making it applicable to a wide range of non‑convex problems where projection onto $\mathcal C$ is expensive but a linear minimization oracle is cheap.

\section*{Preliminaries (Definitions and Lemmas)}
\begin{lemma}[Descent Lemma for $L$‑smooth functions]\label{lem:descent}
For any $x,y\in\mathcal C$,
\[
f(y) \le f(x) + \ip{\nabla f(x)}{y-x} + \frac{L}{2}\norm{y-x}^2 .
\]
\end{lemma}
\begin{proof}
Standard consequence of $L$‑smoothness; see e.g. \cite{Nesterov2004}. \qedhere
\end{proof}

\begin{lemma}[Bound on the distance between iterate and LMO point]\label{lem:diam}
For all $k$, $\norm{s_k-x_k}\le D$.
\end{lemma}
\begin{proof}
Both $s_k$ and $x_k$ belong to $\mathcal C$, and $D$ is the diameter of $\mathcal C$ by Assumption \ref{ass:compact}. \qedhere
\end{proof}

\section*{Main Proof (Step‑by‑Step Derivation)}
We prove Theorem \ref{thm:main}. Throughout we fix an arbitrary iteration $k\ge0$.

\begin{enumerate}[label=[Step \arabic*], leftmargin=*, itemsep=4pt]
\item Apply Lemma \ref{lem:descent} with $x=x_k$ and $y=x_{k+1}=(1-\gamma_k)x_k+\gamma_k s_k$:
\[
f(x_{k+1}) \le f(x_k) + \ip{\nabla f(x_k)}{x_{k+1}-x_k}
                + \frac{L}{2}\norm{x_{k+1}-x_k}^2 . \tag{2}
\]

\item Compute the two inner terms explicitly.
\[
x_{k+1}-x_k = \gamma_k (s_k-x_k). \tag{3}
\]
Hence
\[
\ip{\nabla f(x_k)}{x_{k+1}-x_k}
     = \gamma_k \ip{\nabla f(x_k)}{s_k-x_k}
     = -\gamma_k g_{\mathcal C}(x_k) . \tag{4}
\]

\item Using (3) again,
\[
\norm{x_{k+1}-x_k}^2 = \gamma_k^2 \norm{s_k-x_k}^2
                    \le \gamma_k^2 D^2 \quad\text{by Lemma \ref{lem:diam}}. \tag{5}
\]

\item Substitute (4) and (5) into (2):
\[
f(x_{k+1}) \le f(x_k) - \gamma_k g_{\mathcal C}(x_k) 
                + \frac{L}{2}\gamma_k^2 D^2 . \tag{6}
\]

\item Rearrange (6) to isolate the gap term:
\[
\gamma_k g_{\mathcal C}(x_k) \le f(x_k)-f(x_{k+1}) 
                               + \frac{L}{2}\gamma_k^2 D^2 . \tag{7}
\]

\item Sum inequality (7) over $k=0,\dots,T$:
\[
\sum_{k=0}^{T}\gamma_k g_{\mathcal C}(x_k)
   \le f(x_0)-f(x_{T+1}) + \frac{L D^2}{2}\sum_{k=0}^{T}\gamma_k^2 . \tag{8}
\]

\item Since $f(x_{T+1})\ge f_{\inf}$, we obtain
\[
\sum_{k=0}^{T}\gamma_k g_{\mathcal C}(x_k)
   \le f(x_0)-f_{\inf} + \frac{L D^2}{2}\sum_{k=0}^{T}\gamma_k^2 . \tag{9}
\]

\item By definition of $\tilde g_T$,
\[
g_{\mathcal C}(x_k)\ge \tilde g_T\quad\forall\,k\le T .
\]
Multiplying by $\gamma_k$ and summing yields
\[
\tilde g_T \sum_{k=0}^{T}\gamma_k \le 
      \sum_{k=0}^{T}\gamma_k g_{\mathcal C}(x_k). \tag{10}
\]

\item Combine (9) and (10):
\[
\tilde g_T \sum_{k=0}^{T}\gamma_k 
   \le f(x_0)-f_{\inf} + \frac{L D^2}{2}\sum_{k=0}^{T}\gamma_k^2 . \tag{11}
\]

\item Finally divide by $\sum_{k=0}^{T}\gamma_k$ and multiply the second term by $2$ to obtain the bound announced in the theorem:
\[
\boxed{\;
\tilde g_T \le 
\frac{2\bigl(f(x_0)-f_{\inf}\bigr)}{\sum_{k=0}^{T}\gamma_k}
      + L D^2 \frac{\sum_{k=0}^{T}\gamma_k^2}{\sum_{k=0}^{T}\gamma_k}
\;} . \tag{12}
\]
\end{enumerate}
This completes the proof. \qed

\section*{Rate Derivation (With Constants)}
Choose $\gamma_k = \frac{2}{k+2}$.
Then
\[
\sum_{k=0}^{T}\gamma_k = 2\sum_{k=0}^{T}\frac{1}{k+2}
   \ge 2\int_{1}^{T+2}\frac{1}{t}\,dt = 2\ln(T+2) .
\]
Similarly,
\[
\sum_{k=0}^{T}\gamma_k^2 = 4\sum_{k=0}^{T}\frac{1}{(k+2)^2}
   \le 4\left(1+\int_{1}^{\infty}\frac{1}{t^2}\,dt\right)=8 .
\]
Insert into (12):
\[
\tilde g_T \le 
\frac{2(f(x_0)-f_{\inf})}{2\ln(T+2)} + L D^2\frac{8}{2\ln(T+2)}
          = O\!\left(\frac{1}{\ln T}\right).
\]
If instead we use a constant step size $\gamma_k=\gamma\in(0,1]$, then
\[
\tilde g_T \le \frac{2(f(x_0)-f_{\inf})}{\gamma (T+1)} 
          + L D^2\,\gamma .
\]
Optimizing $\gamma = \Theta(1/\sqrt{T})$ yields the familiar $O(1/\sqrt{T})$ rate.

\section*{Special Cases}
\begin{enumerate}
\item \textbf{Convex $f$.} The Frank–Wolfe gap coincides with the primal–dual gap, and (12) simplifies to the classical $O(1/T)$ bound when $\gamma_k=2/(k+2)$.
\item \textbf{Polytope $\mathcal C$.} The LMO reduces to a linear program over the vertices; the bound (12) remains unchanged but the computational cost per iteration is often much lower than a projection.
\item \textbf{Exact line‑search.} If $\gamma_k$ is chosen by minimizing $f\bigl((1-\gamma)x_k+\gamma s_k\bigr)$, then the term $\frac{L}{2}\gamma_k^2 D^2$ in (6) can be replaced by $0$, improving the constant in (12).
\end{enumerate}

\begin{thebibliography}{9}
\bibitem{Nesterov2004}
Y.~Nesterov, \emph{Introductory Lectures on Convex Optimization: A Basic Course}, Kluwer Academic Publishers, 2004.
\end{thebibliography}

\end{document}